% !TEX root = ../main.tex
\chapter{チーム開発で使うための運用}
\label{chap:team-operations}
\BookIndex{ちーむうんよう}{チーム運用}
\BookIndex{continuous integration}{CI}
\BookIndex{れびゅー}{レビュー}

\begin{chapterabstract}
本章では、Lean 4 で書いた仕様モデルと証明を、個人の学習成果ではなく、チーム開発のレビュー、CI、変更管理に組み込む方法を扱います。
Lean ファイルの置き場所、仕様レビューでの使い方、\texttt{lake build} による継続的検査、壊れた証明の読み方、実装コードと検証モデルの同期、そして「証明済み」と述べてよい範囲の明記を順に整理します。

本章の主張は単純です。
証明は、チームにとって新しい種類のテストではなく、仕様変更に反応するセンサーです。
センサーとして使うには、どの定理がどの仕様項目に対応しているか、どのファイルを誰がレビューするか、証明モデルが実装のどこまでを表しているかを、運用として決めておく必要があります。
\end{chapterabstract}

\begin{learninggoals}
\item Lean ファイルをリポジトリ内のどこに置くかを判断できます。
\item 証明を仕様レビューとコードレビューの補助資料として使えます。
\item CI で \texttt{lake build} を回し、壊れた証明を変更影響のシグナルとして読めます。
\item 実装コードと検証モデルの同期問題を説明し、同期の粒度を決められます。
\item 「証明済み」と主張できる範囲と、証明していない範囲を明文化できます。
\end{learninggoals}

\section{小さな形式化をチームで運用する}

ここまでの章では、仕様の核を Lean で小さく表す方法を学んできました。
たとえば、データ移行の round-trip、API adapter の自然性、リスト移行の件数保存、状態更新の不変条件保存などです。
こうした性質は一人で試すだけでも有用ですが、チームの変更プロセスに組み込むと実務上の価値が高まります。

個人の手元では、Lean ファイルは「動くノート」でも構いません。
しかしチーム開発では、それだけでは足りません。
誰が読めるのでしょうか。
どの仕様に対応しているのでしょうか。
証明が壊れたとき、実装が間違っているのでしょうか、仕様が変わったのでしょうか、それともモデルだけが古くなったのでしょうか。
責任者、仕様との対応、証明が壊れた場合の判断方法を決めずに導入すると、Lean は品質保証の道具ではなく、読めない追加ファイルになります。

Lean の証明をチームで使うとは、「正しいことを一度証明して終わり」ではありません。
変更のたびに、以前の約束がまだ保たれているかを機械に確認してもらう運用を作ることです。

\FigurePlaceholder{fig-ch36-team-loop.png}{0.92}{Lean 検証の運用ループ}{fig:ch36-team-loop}

図~\ref{fig:ch36-team-loop} は、本章で扱う流れを表します。
仕様から証明したい性質を切り出し、Lean ファイルに定理として置きます。
レビューでは、自然言語の仕様と Lean の定理名を対応させます。
CI では \texttt{lake build} を回します。
証明が壊れたら、変更の意味を確認します。
仕様の切り出し、レビュー、CI、変更理由の確認という流れを小さく回すことが、Lean をチーム開発に入れる基本形です。

\section{Lean ファイルの置き場所}

最初に決めるべきことは、Lean ファイルをどこに置くかです。
置き場所には、二つの考え方があります。
一つは、本番コードの近くに置く方法です。
もう一つは、\texttt{verification/} や \texttt{spec/} のような専用ディレクトリに分ける方法です。

本書では、導入初期には専用ディレクトリに分ける方法を勧めます。
Lean のモデルが本番コードそのものではなく、証明しやすいように抽出された検証モデルだからです。
最初から本番コードと同じ階層に置くと、「これは実装なのか、仕様モデルなのか」が曖昧になりやすくなります。

実装コードが \texttt{src/} にあるなら、Lean の検証モデルは \texttt{verification/} や \texttt{lean/} に置くとよいでしょう。
さらに、対象ドメインごとに \texttt{UserMigration.lean}、\texttt{Pricing.lean}、\texttt{ApiAdapter.lean} のように分けます。

例として、次のような構成を考えます。

\begin{listing}[htbp]
\begin{minted}{text}
project-root/
  src/
    user_migration.ts
    pricing.ts
    api_adapter.ts
  docs/
    specs/
      user_migration.md
      pricing.md
  verification/
    lakefile.lean
    lean-toolchain
    Verification/
      UserMigration.lean
      Pricing.lean
      ApiAdapter.lean
      SpecIndex.lean
\end{minted}
\caption{\texttt{verification/} の構成}
\label{lst:ch36-directory-layout}
\end{listing}

この構成では、\texttt{src/} が本番コード、\texttt{docs/specs/} が自然言語仕様、\texttt{verification/} が Lean プロジェクトです。
重要なのは、Lean ファイルを孤立させないことです。
\texttt{SpecIndex.lean} のような索引ファイルを作り、仕様項目と theorem 名を対応させると、レビューで追いやすくなります。

本章では、\textbf{検証モデル}という言葉を次の意味で使います。
検証モデルとは、本番コードの全機能を再実装したものではなく、証明したい性質に関係する型、関数、述語だけを抽出した Lean 上の小さなモデルです。

検証モデルは、本番コードより小さくて構いません。
むしろ、小さくなければレビューできません。
チームに必要なのは、すべてを Lean に移すことではなく、「この仕様項目に対して、この小さなモデルではこの性質が証明されている」と追跡できる状態です。

\section{証明を仕様レビューに使う}

Lean の定理は、仕様レビューで使うと効果が出やすくなります。
コードレビューだけに置くと、実装差分の一部として埋もれるためです。
仕様レビューでは、自然言語の主張を見ながら、それが Lean ではどの theorem に対応するかを確認します。

たとえば、仕様書に次のような文があるとします。

\begin{quote}
ユーザー移行では、旧形式から新形式へ移行したあと、旧形式へ戻すと元のユーザーと一致する。
\end{quote}

Lean では、これを次のような theorem 名に対応させます。

\begin{listing}[htbp]
\begin{minted}{lean4}
namespace Chapter36

structure UserV1 where
  id : Nat
  name : String
  deriving Repr, DecidableEq

structure UserV2 where
  id : Nat
  displayName : String
  deriving Repr, DecidableEq

def migrate (u : UserV1) : UserV2 :=
  { id := u.id, displayName := u.name }

def rollback (u : UserV2) : UserV1 :=
  { id := u.id, name := u.displayName }

theorem userV1_rollback_migrate (u : UserV1) :
    rollback (migrate u) = u := by
  cases u
  rfl

end Chapter36
\end{minted}
\caption{\texttt{userV1\_rollback\_migrate}}
\label{lst:ch36-spec-link}
\end{listing}

定理 \texttt{userV1\_rollback\_migrate} は、引用した仕様を \texttt{UserV1} と \texttt{UserV2} の簡略モデル上でそのまま表しています。
ただし、このモデルの移行は二フィールドをコピーして名前を変えるだけであり、DB I/O、時刻補完、重複 ID、パース失敗は扱いません。
実際のプロジェクトでは、仕様書の項番、定理名、前提、モデル外の事項をレビュー表で対応づけます。

Lean で重要なのは、定理名を人間が読めるレビュー単位にすることです。
定理名を \path{userV1_rollback_migrate} のようにすれば、「UserV1 の rollback と migrate の round-trip についての証明だ」と推測できます。
逆に、\texttt{thm1} や \texttt{aux2} ばかりでは、CI が通ってもチームの知識にはなりません。

仕様レビューでは、次の三つを確認します。

\begin{itemize}
\item 自然言語仕様が、Lean で表せる粒度まで分解されているかを確認します。
\item theorem 名が、仕様項目と対応しているかを確認します。
\item 証明の前提条件が、仕様書にも明示されているかを確認します。
\end{itemize}

自然言語仕様の分解、theorem 名との対応、前提条件の明示がそろって初めて、Lean の証明はレビュー可能な資料になります。

\section{CI で \texttt{lake build} を回す}

Lean の証明をチームで使うなら、CI で自動的に検査します。
手元でだけ \texttt{lake build} を実行している状態では、証明は個人の作業に閉じます。
Pull Request ごとに \texttt{lake build} が走るようにしておくと、仕様モデルや定理が壊れた変更を早く検出できます。

CI で Lean を回す目的は、「Lean ファイルもコンパイル対象にする」ことです。
証明は実行時テストではありませんが、\texttt{lake build} が失敗すれば、少なくとも Lean 上の仕様主張は現在の定義と整合していません。

典型的な GitHub Actions の形は、次のようになります。

\begin{listing}[htbp]
\begin{minted}{text}
name: lean-verification

on:
  pull_request:
  push:
    branches: [ main ]

jobs:
  lake-build:
    runs-on: ubuntu-latest
    defaults:
      run:
        working-directory: verification
    steps:
      - uses: actions/checkout@v4
      - uses: leanprover/lean-action@v1
      - name: Build Lean verification project
        run: lake build
\end{minted}
\caption{\texttt{lean-verification}}
\label{lst:ch36-ci-build}
\end{listing}

この例では、Lean プロジェクトを \texttt{verification/} に置いているため、CI の working directory もそこに合わせています。
本番コード側のテストとは別ジョブにしても、同じワークフローの中で並列に動かしても構いません。

CI で \texttt{lake build} が通ることは、「本番コードが完全に正しい」ことを意味しません。
通っているのは、Lean 上に書いた検証モデルと定理です。
本番コードがそのモデルと対応しているかは、別途レビューと同期ルールで管理します。

導入初期には、CI の対象を広げすぎないほうがよいでしょう。
最初は、壊れると困る中核仕様を一つか二つ選びます。
たとえば、\texttt{UserMigration.lean} の round-trip、\texttt{Pricing.lean} の単純な等式、\texttt{ApiAdapter.lean} の自然性のようなものだけで構いません。
CI が安定してから、対象を増やします。

\section{壊れた証明をレビューシグナルとして読む}

証明が壊れたとき、すぐに「Lean が面倒なことを言っている」と考えてはいけません。
壊れた証明は、変更影響のシグナルです。
ただし、そのシグナルには複数の意味があります。

証明が壊れる原因は、主に次の四つです。

\begin{itemize}
\item 実装に対応するモデル定義が変わり、以前の性質が成り立たなくなりました。
\item 仕様が変わったので、以前の theorem 自体を更新すべきです。
\item 補助定義や名前が変わっただけで、主張はまだ正しい状態です。
\item 証明モデルが本番コードからずれており、モデル側を見直す必要があります。
\end{itemize}

この分類を持っておくと、壊れた証明をレビューで扱いやすくなります。
証明の修正を単なる作業として片付けるのではなく、「何が変わったから壊れたのか」を Pull Request の説明に書きます。

\FigurePlaceholder{fig-ch36-broken-proof-flow.png}{0.92}{証明失敗の判断フロー}{fig:ch36-broken-proof-flow}

図~\ref{fig:ch36-broken-proof-flow} は、壊れた証明を読むときの判断フローです。
最初に、定義変更が意図した仕様変更かどうかを見ます。
次に、theorem の主張を変える必要があるか、証明だけを更新すればよいかを分けます。
最後に、本番コードと検証モデルの対応がまだ妥当かを確認します。

本章では、\textbf{レビューシグナル}を「変更内容について追加確認が必要だと知らせる機械的な兆候」と呼びます。
壊れた証明は、失敗そのものではなく、変更前に成り立っていた仕様主張が現在の定義では確認できなくなったというシグナルです。

重要なのは、証明を無理に直さないことです。
仕様が変わったなら、theorem 名や仕様書も変えます。
仕様は変わっていないのに証明が壊れたなら、実装変更やモデル変更が意味を変えた可能性があります。
その区別をレビューで確認します。

\section{ドメインコードと検証モデルの同期問題}

Lean を実務に持ち込むとき、最も誤解されやすいのが同期問題です。
Lean の検証モデルは、本番コードそのものではありません。
したがって、モデルが正しくても、本番コードが同じ意味で実装されていなければ保証は弱くなります。

検証モデルは地図であり、本番コードは地形です。
地図が正確であれば移動に役立ちますが、地形が変わったのに地図を更新しなければ危険です。
逆に、地図を細かくしすぎると、更新できなくなります。

同期の方法には、いくつかの粒度があります。

\begin{itemize}
\item \textbf{手動同期}: 実装変更時に、担当者が Lean モデルも更新します。
\item \textbf{レビュー同期}: Pull Request テンプレートに「Lean モデルへの影響」を記入します。
\item \textbf{仕様同期}: 仕様書の項番と theorem 名を対応表で管理します。
\item \textbf{生成同期}: 可能な範囲で、型定義や schema から Lean 側の一部を生成します。
\end{itemize}

本書の範囲では、手動同期とレビュー同期から始めます。
生成同期は強力ですが、導入コストが高い方法です。
最初から自動生成を目指すと、Lean の導入目的が「証明」ではなく「ツールチェーン整備」にすり替わりやすくなります。

次の Lean コードは、同期対象を小さな記録として表す例です。

\begin{listing}[htbp]
\begin{minted}{lean4}
namespace Chapter36

structure ModelLink where
  specId : Nat
  theoremName : String
  modelFile : String
  implementationFile : String
  assumption : String


def userMigrationLink : ModelLink :=
  { specId := 2204
    theoremName := "userV1_rollback_migrate"
    modelFile := "Verification/UserMigration.lean"
    implementationFile := "src/user_migration.ts"
    assumption := "Lean model ignores database IO and timestamps" }

#eval userMigrationLink.theoremName  -- => "userV1_rollback_migrate"

end Chapter36
\end{minted}
\caption{\texttt{ModelLink}・\texttt{userMigrationLink}}
\label{lst:ch36-model-link}
\end{listing}

このコードは証明ではありません。
しかし、チーム運用ではこのようなメタデータも重要です。
証明の theorem 名だけが残っていても、どの実装ファイルと対応しているかが分からなければ、レビューで使いにくくなります。
Lean の中に書いても、Markdown や YAML の対応表として管理しても構いません。

\section{「証明済み」と言ってよい範囲を明記する}

Lean を導入すると、強い言葉を使いたくなります。
「この機能は証明済みです」「この移行は安全です」という言い方です。
しかし、その言い方は危険です。
正確には、「この Lean モデルにおいて、この前提のもとで、この theorem が検査されています」と述べるべきです。

「証明済み」という言葉は、必ず範囲と前提を添えて使います。
Lean が証明したのは、Lean に書いた命題です。
本番環境、外部 API、DB、時刻、並行性、パフォーマンス、権限設定まで自動的に保証されるわけではありません。

チームで使う場合、証明済み範囲の記述には次の五項目を含めるとよいでしょう。

\begin{itemize}
\item \textbf{対象}: どの仕様項目、どの関数、どのデータ移行でしょうか。
\item \textbf{モデル}: Lean でどの型または関数として表したのでしょうか。
\item \textbf{定理}: どの theorem が通っているのでしょうか。
\item \textbf{前提}: どの条件を仮定しているのでしょうか。
\item \textbf{範囲外}: 何を証明していないのでしょうか。
\end{itemize}

たとえば、次のように書きます。

\begin{quote}
\textbf{証明済み範囲}: \texttt{Verification/UserMigration.lean} の \texttt{userV1\_rollback\_migrate} により、Lean 上の \texttt{UserV1} と \texttt{UserV2} モデルについて、\texttt{rollback (migrate u) = u} を検査している。ただし、DB からの読み出し、時刻フィールドの自動補完、JSON parser の実装、実データの全件移行はこの証明の範囲外である。
\end{quote}

この書き方なら、証明の価値を過小評価せず、同時に過大評価もしません。
チームにとって重要なのは、証明が何を保証し、何を保証していないかを共有することです。

\section{導入初期の運用パターン}

最後に、導入初期の現実的なパターンをまとめます。
最初から全仕様を Lean に移す必要はありません。
むしろ、壊れると被害が大きく、かつ小さくモデル化できる箇所から始めます。

\subsection{パターン1: マイグレーションの round-trip だけを守る}

最初の題材として最も分かりやすいのは、スキーマ移行です。
\texttt{migrate} と \texttt{rollback} を Lean で小さく書き、round-trip を証明します。
CI でこの theorem を回すだけでも、移行仕様の変更に強いセンサーになります。

\subsection{パターン2: 料金計算のリファクタリングだけを守る}

料金計算や集計ロジックでは、実装全体を Lean に移すのではなく、丸めを含まない中核式や、順序を入れ替えてよい部分だけをモデル化します。
証明できる核を決め、範囲外を明記します。

\subsection{パターン3: API adapter の可換性だけを守る}

レスポンス wrapper の変更では、adapter が業務ロジックと干渉しないことを可換性として証明します。
これは、自然変換の章で扱った考え方を、チームレビューに移す例です。

\section{本章のまとめ}

Lean ファイルは、導入初期には \texttt{verification/} や \texttt{lean/} のような専用ディレクトリに置くと、実装と検証モデルの役割を分けやすくなります。

theorem 名と仕様項目を対応させることで、証明は仕様レビューの資料になります。

CI で \texttt{lake build} を回すと、仕様モデルや定理の破壊を Pull Request の段階で検出できます。

壊れた証明は失敗ではなく、仕様変更、実装変更、モデル更新のどれが起きたかを確認するレビューシグナルです。

「証明済み」と述べるときは、対象、モデル、定理、前提、範囲外を必ず明記します。
